\chapter{Przykłady z równań różniczkowych}

Zachowujemy kilka przykładów pokazujących, jak automorfizmy współgrają z pochodną. Nie są one przesłankami dowodów z rozdziału o całkach. Nie rozwijamy tu ogólnej teorii liniowych grup algebraicznych.

\section{Automorfizmy zachowujące pochodną}

Dla rozszerzenia różniczkowego $L/K$ rozważamy grupę automorfizmów ciała $L$, które ustalają $K$ i spełniają
\[
D\sigma(u)=\sigma(Du)\qquad(u\in L).
\]
Ten dodatkowy warunek wymusza zachowanie równań różniczkowych. W trzech następnych przykładach $K=\C(x)$, $D=\dd/\dd x$, a ciało stałych pozostaje równe $\C$.

\begin{przyklad}[Równanie wykładnicze]
Dla $y'=y$ wybierzmy $t=e^x$ i $L=\C(x,t)$. Jeżeli $\sigma$ jest automorfizmem różniczkowym nad $K$, to
\[
D\left(\frac{\sigma(t)}{t}\right)
=\frac{D\sigma(t)\cdot t-\sigma(t)Dt}{t^2}=0.
\]
Zatem $\sigma(t)=ct$ dla pewnego $c\in\C^\times$. Odwrotnie, każdy taki wybór daje automorfizm $t\mapsto ct$, ponieważ $t$ jest przestępny; jego odwrotnością jest $t\mapsto c^{-1}t$. Sprawdzenie na generatorach pokazuje, że komutuje z $D$. Grupa jest więc dokładnie grupą multiplikatywną $\C^\times$.
\end{przyklad}

\begin{przyklad}[Równanie logarytmiczne]
Dla $y'=1/x$ weźmy lokalną gałąź $\ell=\log x$ i $L=\C(x,\ell)$. Element $\ell$ jest przestępny: gdyby był algebraiczny stopnia $d$, to $b=\Tr(\ell)/d\in\C(x)$ spełniałoby $b'=1/x$. Jest to niemożliwe na mocy lematu 9.1 z $h=0$. Ten sam brak wymiernej funkcji pierwotnej i lemat 8.5 zapewniają brak nowych stałych.

Dla automorfizmu różniczkowego mamy $D(\sigma(\ell)-\ell)=0$, więc $\sigma(\ell)=\ell+c$, $c\in\C$. Przestępność pokazuje, że wszystkie takie przesunięcia są automorfizmami; ich składanie dodaje parametry. Otrzymujemy grupę addytywną $(\C,+)$. Jest to przykład równania niejednorodnego, lecz definicja grupy automorfizmów nadal ma sens.
\end{przyklad}

\begin{przyklad}[Oscylator]
Równanie $y''+y=0$ ma rozwiązania $z=e^{ix}$ i $z^{-1}=e^{-ix}$. Są liniowo niezależne nad $\C$, ponieważ
\[
\det\begin{pmatrix}z&z^{-1}\\iz&-iz^{-1}\end{pmatrix}=-2i\ne0.
\]
Ciało generowane przez nie i ich pochodne to $L=\C(x,z)$. Tak jak w przykładzie wykładniczym wszystkie automorfizmy mają postać $z\mapsto cz$, $c\in\C^\times$. Na drugie rozwiązanie działają jako $z^{-1}\mapsto c^{-1}z^{-1}$. W tej bazie grupa ma więc macierze
\[
\begin{pmatrix}c&0\\0&c^{-1}\end{pmatrix},\qquad c\in\C^\times.
\]
Nie możemy niezależnie skalować obu rozwiązań, bo ich iloczyn wynosi $1$. To ta sama zależność, którą uwzględnialiśmy przy całce sinusowej.
\end{przyklad}

\section{Nieelementarne nie znaczy nieliouvillowskie}

Rozszerzenie liouvillowskie dopuszcza, obok kroków algebraicznych, kroki $D\theta\in F$ oraz $D\theta/\theta\in F$ nad poprzednim ciałem $F$, bez nowych stałych. Są to dowolne funkcje pierwotne i eksponenty funkcji pierwotnych, nie tylko logarytmy i eksponenty elementów $F$. Klasa ta jest szersza od klasy rozszerzeń elementarnych.

Na przykład $F(x)=\int_0^x e^{-t^2}\dd t$ nie jest elementarne, ale jest liouvillowskie z samej definicji: najpierw dołączamy $e^{-x^2}$, a potem jego funkcję pierwotną. Ponadto
\[
F'=e^{-x^2},\qquad F''=-2xe^{-x^2},\qquad F''+2xF'=0.
\]
Nieelementarna funkcja może więc rozwiązywać bardzo proste liniowe równanie różniczkowe. Nie ma w tym sprzeczności z twierdzeniem Liouville'a.

\subsection*{Równanie Airy'ego i równanie Riccatiego}

Dla kontrastu rozważmy $y''-xy=0$ nad $K=\C(x)$. Pokażemy, dlaczego nie ma ono niezerowych rozwiązań liouvillowskich i dlaczego jego różniczkowa grupa Galois to $\mathrm{SL}_2(\C)$. Sam rachunek nie wymaga algorytmu Kovacica, ale użyjemy jednego klasycznego kryterium z teorii Picarda--Vessiota. Nie jest ono konsekwencją zwykłego twierdzenia Galois z rozdziału 3.

\begin{lemat}[Algebraiczne rozwiązania równania Riccatiego]
Równanie
\begin{equation}\label{eq:airy-riccati}
r'+r^2=x
\end{equation}
nie ma rozwiązań algebraicznych nad $\C(x)$.
\end{lemat}
\begin{proof}
Załóżmy, że $r$ jest algebraiczne. Każda jego lokalna gałąź w pobliżu skończonego punktu $x=a$ ma rozwinięcie Puiseux, a więc w punkcie nad $a$ ma granicę w $\C\cup\{\infty\}$.

Jeżeli granica wynosi $c\in\C$, to równanie $r'=x-r^2$ ma analityczną prawą stronę. Twierdzenie o istnieniu i jednoznaczności dla holomorficznego równania pierwszego rzędu daje jedną holomorficzną gałąź z warunkiem $r(a)=c$. Nasza gałąź musi się z nią pokrywać: można to zobaczyć także przez podstawienie szeregu Puiseux do równania, które kolejno wyznacza współczynniki zwykłego szeregu Taylora i wyklucza wykładniki ułamkowe.

Jeżeli $r$ ma biegun, rozważamy $u=1/r$. Z równania wynika $u'=1-xu^2$; tu również prawa strona jest holomorficzna i warunek $u(a)=0$ wyznacza jednoznaczną holomorficzną gałąź. Tak więc $r$ może mieć biegun, lecz nie ma rozgałęzienia nad żadnym skończonym punktem $a$.

W rezultacie skończone przykrycie płaszczyzny $\C$ określone przez wielomian minimalny $r$ jest nierozgałęzione. Ponieważ $\C$ jest po prostu spójna, każde spójne skończone nierozgałęzione przykrycie $\C$ ma jeden arkusz. Stąd $[\C(x,r):\C(x)]=1$, czyli $r\in\C(x)$.

Nie istnieje jednak rozwiązanie wymierne. Niech $d$ będzie stopniem części wielomianowej funkcji wymiernej $r$. Jeśli $d\ge1$, to w nieskończoności $r^2$ ma biegun rzędu $2d$, podczas gdy $r'$ ma co najwyżej rząd $d-1$; lewa strona \eqref{eq:airy-riccati} ma zatem rząd bieguna $2d$, a prawa rząd $1$. Jeśli $r$ jest regularna w nieskończoności, to $r^2+r'$ także jest tam regularna, natomiast $x$ ma biegun. W obu przypadkach otrzymujemy sprzeczność.
\end{proof}

Potrzebujemy teraz następującej standardowej konsekwencji teorii Picarda--Vessiota i twierdzenia Liego--Kołczyna: dla równania $y''=q(x)y$, $q\in\C(x)$, istnieje niezerowe rozwiązanie liouvillowskie wtedy i tylko wtedy, gdy równanie Riccatiego $r'+r^2=q$ ma rozwiązanie algebraiczne nad $\C(x)$. Wyjaśnijmy rolę tego kryterium. Z rozwiązania algebraicznego $r$ otrzymujemy $y=\exp(\int r\,\dd x)$ w rozszerzeniu liouvillowskim. W drugą stronę, jeśli jedno rozwiązanie jest liouvillowskie, to drugie można otrzymać przez redukcję rzędu $y_2=y_1\int y_1^{-2}\dd x$; grupa Picarda--Vessiota ma wtedy rozwiązalną składową spójną. Twierdzenie Liego--Kołczyna zapewnia jej niezmienniczą prostą w dwuwymiarowej przestrzeni rozwiązań. Iloraz $r=y'/y$ odpowiadającego wektora jest ustalany przez tę składową, a więc jest algebraiczny nad ciałem bazowym (grupa ma skończenie wiele składowych). Ten ostatni krok korzysta z podstawowego twierdzenia różniczkowej teorii Galois.

\begin{twierdzenie}[Równanie Airy'ego]
Równanie $y''-xy=0$ nie ma niezerowych rozwiązań liouvillowskich nad $\C(x)$. Jego różniczkowa grupa Galois nad $\C(x)$ jest równa $\mathrm{SL}_2(\C)$.
\end{twierdzenie}
\begin{proof}
Pierwsza teza wynika z kryterium Riccatiego oraz poprzedniego lematu.

Niech $G$ będzie grupą Picarda--Vessiota równania Airy'ego. Wrońskian dwóch rozwiązań spełnia $W'=0$, bo w równaniu nie ma wyrazu z $y'$. Wybierając bazę o wrońskianie równym $1$, otrzymujemy $G\subseteq\mathrm{SL}_2(\C)$. Z braku rozwiązań liouvillowskich wynika, że składowa spójna $G^\circ$ nie jest rozwiązalna. Tymczasem każda właściwa spójna algebraiczna podgrupa $\mathrm{SL}_2(\C)$ jest rozwiązalna: jej algebra Liego jest właściwą podalgebrą trójwymiarowej $\mathfrak{sl}_2(\C)$, ma więc wymiar co najwyżej $2$; każda algebra Liego wymiaru co najwyżej $2$ jest rozwiązalna. Stąd $G^\circ=\mathrm{SL}_2(\C)$ i tym bardziej $G=\mathrm{SL}_2(\C)$.
\end{proof}

\begin{uwaga}
Dowód nieelementarności w rozdziale 9 był samodzielny po udowodnieniu twierdzenia Liouville'a. Tutaj odwołaliśmy się jawnie do innego, głębszego wyniku: kryterium Picarda--Vessiota dla rozwiązań liouvillowskich. Elementarne sprawdzenie, że równanie Riccatiego nie ma rozwiązań algebraicznych, nie zastępuje dowodu tego ogólnego kryterium.
\end{uwaga}

\section{Krótka uwaga o algorytmach}

Redukcja do równań takich jak $r'+hr=s$ jest jednym z narzędzi algorytmu Rischa do całkowania elementarnego. Pełna procedura wymaga efektywnego opisu ciał i rozwiązywania odpowiednich problemów w ciele stałych; sam dowód Liouville'a nie jest jeszcze implementacją algorytmu. Inny problem, istnienie niezerowych rozwiązań liouvillowskich liniowych jednorodnych równań drugiego rzędu o wymiernych współczynnikach, jest przedmiotem algorytmu Kovacica. Nie rozwijamy ani nie dowodzimy poprawności tych procedur.

